\subsection{Why the cubic detector alone gives no pure pentagon obstruction}

\begin{statement}{D23.H11}{Proved}{Only two pure five-point maps extend the marked cubic pair in the same target}
\end{statement}
Let $E_0=\mathbb F_3^2\rtimes SL_2(\mathbb F_3)$, with
the exact marked generators of P10
\[
 X=\left(0,\begin{pmatrix}2&1\\0&2\end{pmatrix}\right),
 \qquad
 Y=\left(\binom10,\begin{pmatrix}2&0\\2&2\end{pmatrix}\right).
\]
There are exactly two homomorphisms
\[
 q:K(0,5)\longrightarrow E_0,
 \qquad q(x_{12})=X,\quad q(x_{23})=Y.
\]
On the cyclically ordered generators
$(x_{12},x_{23},x_{34},x_{45},x_{51})$, their values are
\[
                    (X,Y,1,1,Y),\qquad (X,Y,X,1,1).
                                                        \tag{D23.H11.1}
\]
For either map the image of the full-GT pentagon for a
commutator word $f$ is exactly the two-cycle
$f(Y,X)f(X,Y)=1$. Thus no additional necessary condition
on the marked cubic value can come from a pure five-point
homomorphism into this same target with the prescribed
first pair.

\begin{proof}
Write the five generator images as $(a,b,c,d,e)$.
The presentation of Lochak--Schneps (1997), printed
p.~578, consists of
\[
 [a,c]=[b,d]=[c,e]=[d,a]=[e,b]=1
\]
and
\[
 e b^{-1}a c^{-1}b d^{-1}c e^{-1}d a^{-1}=1.
                                                        \tag{D23.H11.2}
\]
Here $[u,v]=uvu^{-1}v^{-1}$. The source is a presentation
of $K(0,5)$, so satisfying these six relations is sufficient
as well as necessary. Its marking is independently checked
in the certificate by substituting
\[
 \sigma_1^2,\quad\sigma_2^2,\quad\sigma_3^2,
 \quad(\sigma_1\sigma_2\sigma_1)^2,
 \quad(\sigma_2\sigma_3\sigma_2)^2
\]
and using the faithful Artin action of $B_4$ on a free group.

Put $a=X$ and $b=Y$. Their centralizers in $E_0$ are
$\langle X\rangle$ and $\langle Y\rangle$, respectively,
both cyclic of order six. To verify this directly, the
linear part of each has $I-A$ invertible. An affine map
centralizing it must fix its unique fixed point, reducing
the question to the linear centralizer. For $X$ this is
\[
 \left\{\begin{pmatrix}\epsilon&t\\0&\epsilon\end{pmatrix}:
        \epsilon=\pm1,\ t\in\mathbb F_3\right\},
\]
of order six; the lower triangular calculation for $Y$
is identical. Each centralizer contains the corresponding
order-six element, giving the assertions.
Since $X,Y$ generate the center-free group $E_0$ by P10,
their common centralizer is trivial. The commutation
relations therefore imply
\[
            d=1,\qquad c=X^i,\qquad e=Y^j,
             \quad i,j\in\mathbb Z/6.
                                                        \tag{D23.H11.3}
\]

No nonidentity power of $X$ commutes with a nonidentity
power of $Y$. Here is a direct calculation, so that this
fact does not depend only on the finite census.
Write $N=E_{12}$ and $L=E_{21}$ over $\mathbb F_3$.
The linear parts of $X^i,Y^j$ are
\[
             (-1)^i(I-iN),\qquad(-1)^j(I+jL).
\]
If they commute, $ij(NL-LN)=0$, so $i$ or $j$ is
divisible by three. For $i,j\in\{1,\ldots,5\}$,
this means $i=3$ or $j=3$.
The translation parts of $Y,Y^2,Y^3,Y^4,Y^5$ are,
respectively,
\[
                  (1,0),(0,2),(1,1),(0,1),(1,2).
\]
They are all nonzero, whereas $X^3=(0,-I)$.
Commutation with $X^3$ would negate and fix this translation,
which is impossible in characteristic three. If $j=3$
and $i\ne3$, the translation $(1,1)$ would have to be
fixed by the linear part of $X^i$. For odd $i$ that matrix
has no nonzero fixed vector; for $i=2,4$ its fixed vectors
have second coordinate zero. This too is impossible.
Hence $[c,e]=1$ forces $c=1$ or $e=1$.

If $c=1$, equation (D23.H11.2) says that $b^{-1}e$
centralizes $a$. It already centralizes $b$, so it is one
and $e=b$. If $e=1$, the same equation says that
$ac^{-1}$ centralizes $b$; it also centralizes $a$, so
$c=a$. These are exactly (D23.H11.1), and direct substitution
shows that both tuples satisfy every defining relation.

The pentagon in the source convention is
\[
          f(c,d)f(e,a)f(b,c)f(d,e)f(a,b)=1.
\]
A commutator word satisfies $f(t,1)=f(1,t)=1$.
Substitution of either tuple in (D23.H11.1) therefore leaves
exactly $f(Y,X)f(X,Y)=1$, as claimed.
\end{proof}

\paragraph{Certificate and scope.}
The exact replay is
\path{code/affine216_pure_pentagon_probe.py}, with record
\path{records/affine216_pure_pentagon_probe.json}.
It checks the source word in the faithful Artin action,
the two centralizers, all thirty-six possible power pairs,
the two surviving homomorphisms, and the transported
known values. The result concerns the same finite target
and exact cubic marking. It does not dispose of larger
targets, nor of the Hurwitz constructions already used
to obtain C34. The whole proof and exact certificate have passed independent review.
